\section{Визуализатор}

Для анализа и отладки результатов укладки разработан интерактивный 3D-визуализатор на Python с использованием библиотеки Plotly.

% TODO: Добавить скриншот визуализатора с уложенной паллетой (figures/visualizer_example.png)
% Показать: 3D-вид паллеты с разноцветными коробками, панель управления внизу

\subsection{Технологии}

Визуализатор реализован с использованием:
\begin{itemize}
    \item \textbf{Python 3} --- основной язык
    \item \textbf{Plotly} --- библиотека для интерактивных 3D-графиков
    \item \textbf{Pandas} --- работа с CSV-файлами
    \item \textbf{NumPy} --- математические операции
\end{itemize}

Результат --- автономный HTML-файл, который можно открыть в любом браузере без необходимости установки дополнительного ПО.

\subsection{Интерактивные возможности}

\subsubsection{Навигация в 3D-пространстве}
\begin{itemize}
    \item Вращение модели мышью (drag)
    \item Масштабирование (scroll)
    \item Перемещение (shift+drag)
    \item Сброс вида (double-click)
\end{itemize}

\subsubsection{Пошаговый просмотр укладки}

Ключевая функция --- возможность просмотреть процесс укладки по шагам:
\begin{itemize}
    \item Клавиши $\leftarrow$ / A --- убрать последнюю коробку
    \item Клавиши $\rightarrow$ / D --- добавить следующую коробку
    \item Кнопки ``Start (0)'' и ``End (All)'' --- перейти в начало/конец
    \item Поле ввода ``Go to:'' --- перейти к конкретному шагу
\end{itemize}

Это позволяет анализировать, как алгоритм принимает решения на каждом шаге, и выявлять потенциальные проблемы.

% TODO: Добавить скриншот пошагового просмотра (figures/visualizer_steps.png)
% Показать: паллету с частично уложенными коробками и панель управления

\subsubsection{Режимы отображения}

\textbf{Режим граней (meshes):}
\begin{itemize}
    \item Коробки отображаются как полупрозрачные 3D-объекты
    \item Регулируемая прозрачность (\texttt{--opacity})
    \item Цветовая палитра по SKU или единый цвет
\end{itemize}

\textbf{Режим рёбер (edges):}
\begin{itemize}
    \item \texttt{sku} --- рёбра окрашены по SKU, можно выбирать типы в легенде
    \item \texttt{single} --- все рёбра одного цвета
    \item \texttt{none} --- рёбра не отображаются
\end{itemize}

\subsubsection{Hover-информация}

При наведении на коробку отображается:
\begin{itemize}
    \item SKU (артикул товара)
    \item Координаты: $x$, $y$, $z$ (мин и макс)
    \item Размеры коробки
\end{itemize}

\subsection{Параметры командной строки}

\begin{table}[ht]
\centering
\caption{Основные параметры визуализатора}
\label{tab:visualizer_params}
\begin{tabular}{|l|l|p{6cm}|}
\hline
\textbf{Параметр} & \textbf{По умолч.} & \textbf{Описание} \\
\hline
\texttt{csv} & --- & Путь к CSV с результатами укладки \\
\texttt{--mode} & meshes & meshes или wireframe \\
\texttt{--opacity} & 0.35 & Прозрачность граней (0--1) \\
\texttt{--color-by} & none & Цвет граней: none или sku \\
\texttt{--edges} & sku & Режим рёбер: sku, single, none \\
\texttt{--edge-color} & \#000000 & Цвет для режима single \\
\texttt{--edge-width} & 2.0 & Толщина рёбер \\
\texttt{--pallet-length} & --- & Длина паллеты (X) \\
\texttt{--pallet-width} & --- & Ширина паллеты (Y) \\
\texttt{--initial-count} & 1 & Сколько коробок показать изначально \\
\texttt{--out} & auto & Путь к выходному HTML \\
\hline
\end{tabular}
\end{table}

\subsection{Примеры использования}

\textbf{Базовый запуск:}
\begin{verbatim}
python3 src/scripts/visualizer.py answers/261.csv \
    --pallet-length 1200 --pallet-width 800
\end{verbatim}

\textbf{С раскраской по SKU и полупрозрачностью:}
\begin{verbatim}
python3 src/scripts/visualizer.py answers/261.csv \
    --pallet-length 1200 --pallet-width 800 \
    --color-by sku --opacity 0.4 --edges sku
\end{verbatim}

\textbf{Показать сразу всю паллету:}
\begin{verbatim}
python3 src/scripts/visualizer.py answers/261.csv \
    --pallet-length 1200 --pallet-width 800 \
    --initial-count 1000000 --opacity 0.5
\end{verbatim}

\subsection{Формат входных данных}

Визуализатор читает CSV-файл с обязательными колонками:

\begin{verbatim}
SKU,x,y,z,X,Y,Z
900001,0,0,0,600,400,200
900001,600,0,0,1200,400,200
900002,0,400,0,200,800,200
...
\end{verbatim}

где $(x, y, z)$ --- координаты нижнего угла коробки, $(X, Y, Z)$ --- координаты верхнего угла.

% TODO: Добавить скриншот с раскраской по SKU (figures/visualizer_sku_colors.png)
% Показать: паллету с коробками разных цветов и легенду справа
